\documentclass{article}
\usepackage{graphicx} % Required for inserting images
\usepackage{amsmath} % Required for inserting images
\usepackage{amssymb}
\title{Physics Informed Priors for MLIPs with UQ}
\author{George Simmons}
\date{June 2026}

\begin{document}
\maketitle
\vspace{2pc}
\section{Meta / abstract points}
\begin{enumerate}
    \item To propagate uncertainties up scales we need to capture MLIP errors at the level of MLIP parameters
    \item However, the resultant model ensemble can be highly unphysical.
    \item Here, we implement physics informed priors for linear MLIPs through inequality constraints
    \item We provide an open source implementation within the POPS UQ scheme
    \item For simple models, the constraints notably improve phonon stability 
    \item The final scheme provides robust max/min bounds while retaining physically meaningful model parameters. 
\end{enumerate}
\section{Main points}
\begin{enumerate}
    \item To communicate that addressing holes in the potential energy surface is traditionally addressed by evaluating more DFT (via hyper-active learning, fine-tuning, etc.) and can be done less expensively
    \item DFT stresses are not well converged so elastic properties are often incorrectly predicted by DFT
    \item Non-repulsive core interactions for potentials are often observed, using our priors allows us to access high pressure conditions
    \item Priors can be formulated based on experimental data
    \item Inclusion of UQ via POPS and analysis of how this affects coverage of training data/test data
\end{enumerate}

\section{Method}
\begin{equation}
        \text{Minimize}\quad  \left\| Xc-y \right\|_2^2 + \left\| \Gamma c \right\|_2^2
\end{equation}

Enforcing linear constraints during training can be achieved by utilising the quadratic programming with linear constraints framework. This can be described mathematically by
\begin{equation}
\begin{split}
    & \text{minimize} \quad \frac{1}{2}c^T P c + q^Tc \\
    & \text{subject to} \quad l \leq Ac \leq u,
\end{split}  
\end{equation}
where $c \in \mathbb{R}^{d}$ is the model parameters that we are minimizing. The matrix $P \in \mathbb{S}_+^d$ is the regularised Gram matrix given by $P := B^TB + \Gamma^T\Gamma$, and $q \in \mathbb{R}^{d}$ is given by $q := B^T y$.
\section{Results}
The main result we present in this work is the enforcement of thermodynamic stability of our POPS ensemble members under phononic displacements.
Fig.~\ref{fig:phonons before and after} displays the improvement of lattice stability at 0K of our ensemble. Qualitatively, this is enforcing that our potential predicts that our crystal geometry remains face-centered cubic (FCC) at 0K. 

However, in practice we care about lattice stability at non-zero temperatures below melting temperature.
Since the only change to Eq.~\ref{eq:linear hessian} is the geometry in the Hessian, the constraint is still linear in parameters and we can enforce positivity easily.

\begin{figure}
    \centering
    \includegraphics[width=\linewidth]{fcc_stability_300K.pdf}
    \caption{a) Phonon bands at lattice constant predicted by 300K}
    \label{fig:constrained fcc stability}
\end{figure}

\begin{figure}
    \centering
    \includegraphics[width=\linewidth]{fcc_structure_300K.pdf}
    \caption{Caption}
    \label{fig:constrained fcc stability}
\end{figure}





\begin{figure}
    \centering
    \includegraphics[width=0.5\linewidth]{bands_naive.pdf}%
    \includegraphics[width=0.5\linewidth]{bands_constrained.pdf}
    \caption{a) Ridge regression solution of Al ACE model with POPS hypercube ensemble. b) Constrained ridge regression solution with constrained POPS hypercube ensemble. The constraints are that the phonon energies along the band-path are greater than 0 eV.}
    \label{fig:phonons before and after}
\end{figure}

\begin{figure}
    \centering
    \includegraphics[width=0.5\linewidth]{energy_parity.pdf}%
    \includegraphics[width=0.5\linewidth]{force_parity.pdf}
    \caption{Al 12, 4, 6, 2}
    \label{fig:model parities}
\end{figure}
\begin{figure}
    \centering
    \includegraphics[width=0.5\linewidth]{energy_calibration.pdf}%
    \includegraphics[width=0.5\linewidth]{force_calibration.pdf}
    \caption{Al 12, 4, 6, 2}
    \label{fig:phonons parities}
\end{figure}


\begin{figure}
    \centering
    \includegraphics[width=1\linewidth]{eos_with_hist.png}
    \caption{Bulk Tungsten cell energies for different lattice constants, with histogram detailing the distribution of pair distances in the training data set.}
    \label{fig:placeholder}
\end{figure}
\end{document}
